% TOP_PROBLEM: {"id": 30002922, "title": "Full-Measure Families of Ergodic Irrational Polygons", "category_id": 11, "category": {"id": 11, "name": "dynamical_systems", "display_name": "Dynamical Systems", "description": "Problems about long-term behavior of deterministic systems, Hamiltonian dynamics, and periodic orbits.", "slug": "dynamical-systems", "order_index": 11, "created_at": "2026-07-31T15:26:25.671Z"}, "status": "open"}
% FIRST_POSTED: 2026-10-06
% ATTEMPT_STATUS: unresolved
% OPENAI_CLAIM: {"schema_version":1,"source_repo":"https://github.com/openai/math","source_commit":"adc7f1241b42e322a6451854ab7e4b4c146bf78a","release_date":"2026-10-06","match":"partial","resolution":"partial","independently_reviewed":false,"second_pass":"n/a","families":[{"family":"150","title":"Weak mixing of triangular billiards with an irrational angle","subject":"Dynamical systems and ergodic theory","manuscripts":[{"title":"Ergodicity of triangular billiards with an irrational angle","pdf_path":"preprints/Ergodicity-of-triangular-billiards-with-an-irrational-angle-September-25-2026/Ergodicity-of-triangular-billiards-with-an-irrational-angle-September-25-2026.pdf","readme_path":"preprints/Ergodicity-of-triangular-billiards-with-an-irrational-angle-September-25-2026/README.md","date":"2026-09-25","theorem_ref":"Theorem labelled thm:main"},{"title":"Weak mixing of triangular billiards with an irrational angle","pdf_path":"preprints/Weak-mixing-of-triangular-billiards-with-an-irrational-angle-October-5-2026/weak-mixing-triangular-billiards.pdf","readme_path":"preprints/Weak-mixing-of-triangular-billiards-with-an-irrational-angle-October-5-2026/README.md","date":"2026-10-05","theorem_ref":"Theorem labelled thm:main"}],"lean":{"doc_path":"lean/docs/150.md","comparators":["lean/ComparatorChallenges/IrrationalTriangleBilliard.json","lean/ComparatorChallenges/IrrationalTriangleBilliard.lean"],"declarations":[],"covers_main_theorem":false},"reasoning_trace":null}]}
\documentclass[11pt]{article}
\usepackage[margin=25mm]{geometry}
\usepackage{fontspec}
\setmainfont{Latin Modern Roman}
\usepackage{amsmath,amssymb,mathtools}
\usepackage{unicode-math}
\setmathfont{Latin Modern Math}
\usepackage{enumitem}
\usepackage{xurl,hyperref,bookmark}
\hypersetup{colorlinks=true,urlcolor=blue}
\setcounter{secnumdepth}{0}
\setlength{\parindent}{0pt}
\setlength{\parskip}{5pt}
\setlength{\emergencystretch}{3em}
\begin{document}
\section*{Full-Measure Families of Ergodic Irrational Polygons}

\subsection{OpenAI's claim}
Manuscript 1 (family 150, 2026-09-25): The full unit-speed billiard flow of each nondegenerate triangle with an angle irrational relative to $\pi$ is claimed to be ergodic, without a Diophantine or genericity condition. Manuscript 2 (family 150, 2026-10-05): An irrational angle divided by $\pi$ is claimed to force weak mixing for the full billiard flow of any nondegenerate Euclidean triangle, using normalized Liouville measure.

\subsection{Outline of the argument}
Manuscript 1 (family 150, 2026-09-25): On the flat double, positional smoothing and long cone-avoiding segments bound endpoint gradients of angular Fourier tails. Two successive weak-limit passages show that all coefficient gradients vanish without a Sobolev assumption; irrational cone rotation removes every nonzero angular mode, leaving a constant invariant function.

Manuscript 2 (family 150, 2026-10-05): Local transport estimates force a possible eigenfunction to be a plane wave with only direction-dependent amplitude. Reflections through the irrational vertex make that amplitude constant, and reflection at the third side excludes any nonzero frequency; the spectral criterion then gives weak mixing.

\subsection{Scope relative to this problem}
Manuscript 1 (family 150, 2026-09-25): For triangles (three edges), rational angle triples form a Lebesgue-null set and the theorem gives ergodicity of all remaining triangles; polygons with arbitrary other prescribed edge counts are not handled.

Manuscript 2 (family 150, 2026-10-05): For triangles (three edges), rational angle triples form a Lebesgue-null set and the theorem gives ergodicity of all remaining triangles; polygons with arbitrary other prescribed edge counts are not handled.

\subsection{Formal verification}
For family 150, the scope page documents the following selections. The selection establishes ergodicity and the almost-everywhere flow law for an irrational-angle triangle, including the exceptional-trajectory null set. The weak-mixing companion is not selected. The documented statement covers the main claim of manuscript 1. Its selected declaration also covers the clause used for this problem. Manuscript 2 is not among the papers selected by that page, so the family link does not formalize its main conclusion. It does not establish formal coverage of the complete problem clause. Comparator files: IrrationalTriangleBilliard.lean. Declarations: OAI.TriangularBilliards.irrational\_triangle\_billiard.

This site has not executed the Lean files or checked their dependencies. These descriptions report the scope of the linked declarations, rather than a new verification of the proofs.

\subsection{Status caveat}
OpenAI attributes these results to an internal model. Its release notice cautions that results lacking full formalization can contain mistakes. This site has not independently assessed the proofs; a solved label records the scope of the reported claim.

\subsection{Sources}
\begin{itemize}
\item OpenAI, \emph{Ergodicity of triangular billiards with an irrational angle}, 2026-09-25; citation key \texttt{OAI:Ergodicity-of-triangular-billiards-with-an-irrational-angle-September-25-2026}. \href{https://github.com/openai/math/blob/adc7f1241b42e322a6451854ab7e4b4c146bf78a/preprints/Ergodicity-of-triangular-billiards-with-an-irrational-angle-September-25-2026/Ergodicity-of-triangular-billiards-with-an-irrational-angle-September-25-2026.pdf}{PDF (GitHub)}; \href{https://github.com/openai/math/raw/adc7f1241b42e322a6451854ab7e4b4c146bf78a/preprints/Ergodicity-of-triangular-billiards-with-an-irrational-angle-September-25-2026/Ergodicity-of-triangular-billiards-with-an-irrational-angle-September-25-2026.pdf}{PDF (direct)}; \href{https://github.com/openai/math/blob/adc7f1241b42e322a6451854ab7e4b4c146bf78a/preprints/Ergodicity-of-triangular-billiards-with-an-irrational-angle-September-25-2026/README.md}{Manuscript page and citation}; \href{https://github.com/openai/math/blob/main/preprints/Ergodicity-of-triangular-billiards-with-an-irrational-angle-September-25-2026/Ergodicity-of-triangular-billiards-with-an-irrational-angle-September-25-2026.pdf}{Latest version}.
\item OpenAI, \emph{Weak mixing of triangular billiards with an irrational angle}, 2026-10-05; citation key \texttt{OAI:Weak-mixing-of-triangular-billiards-with-an-irrational-angle-October-5-2026}. \href{https://github.com/openai/math/blob/adc7f1241b42e322a6451854ab7e4b4c146bf78a/preprints/Weak-mixing-of-triangular-billiards-with-an-irrational-angle-October-5-2026/weak-mixing-triangular-billiards.pdf}{PDF (GitHub)}; \href{https://github.com/openai/math/raw/adc7f1241b42e322a6451854ab7e4b4c146bf78a/preprints/Weak-mixing-of-triangular-billiards-with-an-irrational-angle-October-5-2026/weak-mixing-triangular-billiards.pdf}{PDF (direct)}; \href{https://github.com/openai/math/blob/adc7f1241b42e322a6451854ab7e4b4c146bf78a/preprints/Weak-mixing-of-triangular-billiards-with-an-irrational-angle-October-5-2026/README.md}{Manuscript page and citation}; \href{https://github.com/openai/math/blob/main/preprints/Weak-mixing-of-triangular-billiards-with-an-irrational-angle-October-5-2026/weak-mixing-triangular-billiards.pdf}{Latest version}.
\item Family 150: \href{https://github.com/openai/math/blob/adc7f1241b42e322a6451854ab7e4b4c146bf78a/lean/docs/150.md}{Lean scope}.
\item \href{https://github.com/openai/math/blob/adc7f1241b42e322a6451854ab7e4b4c146bf78a/lean/ComparatorChallenges/IrrationalTriangleBilliard.json}{Comparator IrrationalTriangleBilliard.json}.
\item \href{https://github.com/openai/math/blob/adc7f1241b42e322a6451854ab7e4b4c146bf78a/lean/ComparatorChallenges/IrrationalTriangleBilliard.lean}{Comparator IrrationalTriangleBilliard.lean}.
\item \href{https://github.com/openai/math/blob/adc7f1241b42e322a6451854ab7e4b4c146bf78a/CONTENTS.md}{OpenAI catalogue entry}.
\end{itemize}
\end{document}
